\documentclass[12pt, a4paper]{article}

\usepackage[utf8]{inputenc}
\usepackage[T1]{fontenc}
\usepackage[french]{babel}
\usepackage{geometry}
\usepackage{booktabs}
\usepackage{hyperref}
\usepackage{graphicx}
\usepackage{amsmath}
\usepackage{parskip}
\usepackage{titlesec}
\usepackage{fancyhdr}
\usepackage{xcolor}

\geometry{margin=2.5cm}

\pagestyle{fancy}
\fancyhf{}
\rhead{\small Analyse des Réseaux Sociaux}
\lhead{\small Royal Air Maroc — Étude de diffusion}
\cfoot{\thepage}

\titleformat{\section}{\large\bfseries}{}{0em}{}[\titlerule]
\titleformat{\subsection}{\normalsize\bfseries}{}{0em}{}

\title{
    \vspace{-1cm}
    \textbf{Analyse des Réseaux Sociaux} \\[0.4em]
    \large Étude de la diffusion d'une plainte virale \\
    contre Royal Air Maroc (RAM) sur Instagram
}
\author{}
\date{}

% ----------------------------------------------------------------
\begin{document}

\maketitle
\thispagestyle{fancy}

% ----------------------------------------------------------------
\section{Contexte et Problématique}

Une actrice marocaine, personnalité publique disposant d'une large communauté de
\textit{followers} sur Instagram, a été victime d'un incident de surbooking avec la
compagnie Royal Air Maroc (RAM). Face à cette expérience négative, elle a choisi de
partager son mécontentement publiquement via un \textit{Reel} Instagram, déclenchant
une vague de réactions au sein de son réseau.

L'objectif de ce travail est d'étudier la \textbf{dynamique de diffusion} de cette
information au sein du réseau social, en mobilisant deux cadres théoriques :
la \textbf{théorie des réseaux sociaux} et la \textbf{théorie de la diffusion de
l'information}.

% ----------------------------------------------------------------
\section{Méthodologie}

\subsection{Collecte des données}

Les données ont été collectées via un \textit{scraper} développé en Python, utilisant la
bibliothèque \textbf{Playwright} pour intercepter les appels internes de l'API Instagram.
La collecte a ciblé l'ensemble des commentaires et réponses du Reel publié.

\begin{table}[h!]
\centering
\begin{tabular}{@{}ll@{}}
\toprule
\textbf{Indicateur} & \textbf{Valeur} \\
\midrule
Commentaires collectés  & 846 \\
Utilisateurs uniques    & 733 \\
Fils de réponses        & 32 \\
Total des réponses      & 74 \\
Période d'activité      & Novembre 2024 -- Janvier 2026 \\
\bottomrule
\end{tabular}
\caption{Statistiques du corpus collecté}
\end{table}

\subsection{Construction du réseau}

Le réseau a été construit sous forme de \textbf{graphe orienté} (\textit{DiGraph})
avec la bibliothèque \textbf{NetworkX}, en définissant :

\begin{itemize}
    \item Un \textbf{nœud} par utilisateur ayant commenté
    \item Une \textbf{arête de commentaire} : utilisateur $\rightarrow$ post
    \item Une \textbf{arête de réponse} : utilisateur A $\rightarrow$ utilisateur B
          (quand A répond à B)
    \item Une \textbf{arête de mention} : utilisateur A $\rightarrow$ utilisateur B
          (quand A mentionne @B dans son texte)
\end{itemize}

La détection des communautés a été réalisée via l'algorithme de \textbf{Louvain}.

% ----------------------------------------------------------------
\section{Résultats}

\subsection{Structure globale du réseau}

\begin{table}[h!]
\centering
\begin{tabular}{@{}ll@{}}
\toprule
\textbf{Métrique} & \textbf{Valeur} \\
\midrule
Nœuds (utilisateurs)   & 758 \\
Arêtes (interactions)  & 802 \\
Densité                & 0,0014 \\
Composantes connexes   & 1 \\
Nombre de communautés  & 20 \\
\bottomrule
\end{tabular}
\caption{Métriques globales du réseau}
\end{table}

La densité extrêmement faible ($d = 0{,}0014$) indique un réseau \textbf{très peu
interconnecté} : la majorité des utilisateurs n'interagissent pas entre eux et se
contentent de commenter sans engager de dialogue.

\subsection{Leaders d'opinion et centralité}

\begin{table}[h!]
\centering
\begin{tabular}{@{}llll@{}}
\toprule
\textbf{Utilisateur} & \textbf{Degré} & \textbf{Betweenness} & \textbf{In-degree} \\
\midrule
sa\_na\_jabrane    & 0,0132 & 0,0027 & 8 \\
getfitwith\_hala   & 0,0079 & 0,0105 & 5 \\
saadiataznacht     & 0,0079 & 0,0026 & — \\
une\_maman3        & 0,0066 & 0,0053 & 4 \\
\bottomrule
\end{tabular}
\caption{Principaux acteurs du réseau selon les métriques de centralité}
\end{table}

\texttt{sa\_na\_jabrane} est le nœud le plus sollicité (in-degree = 8),
tandis que \texttt{getfitwith\_hala} est le \textbf{pont principal} entre les
communautés (betweenness le plus élevé : 0,0105).

\subsection{Structure des communautés}

\begin{table}[h!]
\centering
\begin{tabular}{@{}lll@{}}
\toprule
\textbf{Communauté} & \textbf{Taille} & \textbf{Membres clés} \\
\midrule
0 (passive)  & 680 & Commentateurs isolés \\
11 (active)  & 20  & sa\_na\_jabrane, saadiataznacht, naddoustyle \\
9            & 7   & farah\_faress\_usa, ouafa\_paffa \\
1            & 6   & oudasa4, drissamkh \\
10           & 6   & getfitwith\_hala, lys\_blanche\_1 \\
\bottomrule
\end{tabular}
\caption{Principales communautés détectées par l'algorithme de Louvain}
\end{table}

% ----------------------------------------------------------------
\section{Analyse Théorique}

\subsection{Théorie des Réseaux Sociaux}

La structure du réseau révèle un modèle \textbf{hub-and-spoke} : l'actrice constitue
le hub central, et les utilisateurs gravitent autour d'elle sans former de liens denses
entre eux. On observe :

\begin{itemize}
    \item Des \textbf{leaders d'opinion}  qui concentrent les échanges et orientent le débat.
    \item Des \textbf{liens faibles dominants}  : la majorité des
          interactions sont uniques et non répétées, limitant la profondeur de la
          diffusion.
    \item Une \textbf{structure communautaire fragmentée} : 90\,\% des utilisateurs
          (680/759) appartiennent à une communauté passive sans interactions réelles.
\end{itemize}



% ----------------------------------------------------------------
\section{Conclusion}

L'analyse des métriques du réseau permet de conclure que la plainte de l'actrice a
généré une \textbf{réaction significative mais superficielle}. L'information a touché
près de 733 utilisateurs uniques sur une période de 14 mois, ce qui témoigne d'une
portée initiale réelle.

Cependant, la faible densité du réseau ($d = 0{,}0014$), le faible taux de réponses
(74 pour 846 commentaires) et l'absence de ponts puissants entre communautés indiquent
que la diffusion ne s'est pas transformée en cascade informationnelle. La compagnie RAM
a subi une \textbf{pression réputationnelle à court terme}, circonscrite à la communauté
de \textit{followers} de l'actrice, sans propagation vers des réseaux plus larges.

Ce travail illustre que la viralité d'un contenu dépend moins de la qualité du message
que de la \textbf{structure du réseau} dans lequel il circule.

% ----------------------------------------------------------------
\begin{thebibliography}{9}

\bibitem{granovetter1973}
Granovetter, M. (1973).
\textit{The Strength of Weak Ties}.
American Journal of Sociology, 78(6), 1360--1380.

\bibitem{rogers1962}
Rogers, E. M. (1962).
\textit{Diffusion of Innovations}.
Free Press, New York.

\bibitem{newman2010}
Newman, M. E. J. (2010).
\textit{Networks: An Introduction}.
Oxford University Press.

\end{thebibliography}

\end{document}
